\chapter{QCD Phase Structure, Plasma, and Dense Matter}
\label{ch:qcd-phase-structure-plasma-dense-matter}

\ConventionsMixed

This chapter records the finite-temperature and finite-density phase
questions of QCD in the same convention as the Yang--Mills chapters:
\[
  S_{\rm YM}
  =
  \frac{1}{4g^2}\int\operatorname{tr}(F_{\mu\nu}F^{\mu\nu}),
  \qquad
  \operatorname{tr}_{\square}(T^aT^b)=\delta^{ab}
\]
for \(SU(N_c)\) unless another invariant bilinear form is explicitly stated.
The chapter is deliberately not a catalogue of phase-diagram lore.  Each
order parameter is tied to a symmetry, a limiting procedure, or a specified
source.  Each weak-coupling statement is separated from the nonperturbative
questions that remain after the weak-coupling scale hierarchy has been
identified.

\section{Thermal QCD Data}

Let \(N_f\) Dirac quarks transform in the fundamental representation of
\(SU(N_c)\).  Put the Euclidean theory on
\[
  M_\beta=S^1_\beta\times \mathbb R^3,
  \qquad
  \beta=T^{-1},
\]
with periodic gauge fields and antiperiodic quark fields around
\(S^1_\beta\).  A baryon chemical potential \(\mu_B\) couples to the
conserved baryon number \(B\), normalized so that a quark has baryon charge
\(1/N_c\).  Equivalently, every quark flavor has quark chemical potential
\[
  \mu_q=\frac{\mu_B}{N_c}.
\]
At finite spatial volume \(V\), a regulator \(a\), and quark masses
\(m=(m_1,\ldots,m_{N_f})\), define
\begin{equation}
  Z_{a,V}(T,\mu_B,m)
  =
  \operatorname{Tr}_{\Hilb_{a,V}}
  \exp\{-\beta(H_{a,V}-\mu_B B_{a,V})\},
  \label{eq:qcd-phase-finite-volume-partition}
\end{equation}
when the Hamiltonian regulator is used.  In a Euclidean lattice regulator
the same object is represented by the corresponding gauge integral with the
temporal links weighted by \(e^{\pm a\mu_q}\) for forward and backward quark
propagation.  The pressure is the thermodynamic limit
\begin{equation}
  p(T,\mu_B,m)
  =
  \lim_{a\to0}\lim_{V\to\infty}
  \frac{T}{V}\log Z_{a,V}(T,\mu_B,m),
  \label{eq:qcd-pressure-definition}
\end{equation}
provided the limit exists along the chosen continuum trajectory.

\begin{definition}[QCD phase datum]
\label{def:qcd-phase-datum}
A QCD phase datum at fixed \(N_c,N_f\) consists of:
\begin{enumerate}
  \item the thermodynamic pressure \(p(T,\mu_B,m)\) as a function on the
  domain where the continuum thermal state is defined;
  \item a list of symmetries that are exact at the chosen masses and chemical
  potentials;
  \item order parameters or source derivatives associated to those exact
  symmetries;
  \item the limiting prescription used to detect spontaneous symmetry
  breaking;
  \item a status label for each nonanalytic locus: theorem, lattice
  continuum-extrapolation claim, controlled weak-coupling claim, or
  conjectural continuation.
\end{enumerate}
The word ``phase'' is used in this chapter only through such a datum.
\end{definition}

\section{Finite-Regulator Analyticity and Thermodynamic Phases}
\label{sec:qcd-phase-finite-regulator-analyticity}

A phase transition is not a property of a single finite lattice integral.  At
finite cutoff and finite volume the partition function is an analytic
function of its sources away from its zeros.  Nonanalyticity can only emerge
from a limiting procedure: thermodynamic volume, continuum cutoff, chiral
mass, or a source-selection limit.  This elementary point is important
because it prevents a phase diagram from being treated as a drawing whose
curves exist before the limiting observables have been defined.

\begin{definition}[Finite lattice thermal gauge datum]
\label{def:qcd-finite-lattice-thermal-gauge-datum}
Fix a finite Euclidean lattice
\[
  \Lambda=(N_\tau,N_s^3,a)
\]
with compact gauge group \(G\), Haar measure \(dU\) on each link, and
antiperiodic temporal boundary condition for quarks.  Let
\(\kappa=(\kappa_\alpha)\) denote finitely many gauge couplings multiplying
bounded continuous class functions \(S_\alpha(U)\), let
\(\zeta=e^{a\mu_q}\) be the quark fugacity, and let
\[
  M_\Lambda(U;m,\zeta)
\]
be the finite fermion matrix in the chosen lattice discretization, with
entries polynomial in the masses \(m_i\) and Laurent-polynomial in \(\zeta\).
For an integer number of unrooted Dirac flavors, the finite-regulator
partition function is
\begin{equation}
  Z_\Lambda(\kappa,m,\zeta)
  =
  \int_{G^{E(\Lambda)}}\prod_{\ell}dU_\ell\,
  \exp\left[-\sum_\alpha \kappa_\alpha S_\alpha(U)\right]
  \prod_{i=1}^{N_f}\det M_{\Lambda,i}(U;m_i,\zeta),
  \label{eq:qcd-finite-lattice-partition-analytic}
\end{equation}
where \(E(\Lambda)\) is the finite link set.  Rooted or otherwise
non-polynomial fermion prescriptions are not included in this datum unless a
separate branch and locality construction has been supplied.
\end{definition}

\begin{proposition}[Finite-regulator analyticity]
\label{prop:qcd-finite-regulator-analyticity}
For the datum of Definition~\ref{def:qcd-finite-lattice-thermal-gauge-datum},
\(Z_\Lambda(\kappa,m,\zeta)\) is entire in the gauge couplings and masses
and Laurent-polynomial in the fugacity \(\zeta\).  Consequently the finite
pressure
\[
  p_\Lambda(\kappa,m,\zeta)
  =
  \frac{T}{V_\Lambda}\log Z_\Lambda(\kappa,m,\zeta)
\]
is holomorphic on every simply connected parameter domain that does not meet
the zero locus of \(Z_\Lambda\), after a branch of the logarithm has been
chosen.
\end{proposition}

\begin{proof}
The lattice has finitely many links, and \(G^{E(\Lambda)}\) is compact.  For
each fixed gauge field \(U\), the determinant factors are polynomials in the
masses and Laurent polynomials in \(\zeta\).  The exponential of the gauge
action has an absolutely and uniformly convergent power series in
\(\kappa\) on compact subsets of complex coupling space, because every
\(S_\alpha(U)\) is bounded on the compact integration domain.  Termwise
integration is justified by uniform convergence on compact parameter sets.
Thus the integral defines an entire function in \(\kappa\) and \(m\), with
only finitely many positive and negative fugacity powers.  Wherever this
function is nonzero on a simply connected domain, a holomorphic logarithm
exists, and the displayed finite pressure is holomorphic there.
\end{proof}

\begin{definition}[Thermodynamic phase transition and Lee--Yang obstruction]
\label{def:qcd-thermodynamic-phase-transition}
Let \(x\) denote a real collection of thermal parameters, masses, and
sources, and let \(p_\Lambda(x)\) be the finite-regulator pressure along a
specified sequence of lattices and continuum trajectories.  A thermodynamic
phase transition at \(x_0\) means that the limiting pressure
\[
  p(x)=\lim_{\Lambda\to\infty}p_\Lambda(x)
\]
exists in a real neighborhood of \(x_0\) but is not real analytic at
\(x_0\), or that the limiting local Gibbs/KMS state is not unique after the
source-selection prescription has been removed.  The finite-regulator
Lee--Yang locus is
\[
  \mathcal Z_\Lambda
  =
  \{z\in\mathcal U_{\mathbb C}\mid Z_\Lambda(z)=0\}.
\]
\end{definition}

\begin{proposition}[Zero-free neighborhoods exclude pressure singularities]
\label{prop:qcd-lee-yang-zero-free-analytic-limit}
Let \(\mathcal U_{\mathbb C}\) be a complex neighborhood of a real point
\(x_0\).  Suppose that, for all sufficiently large \(\Lambda\),
\(Z_\Lambda\) has no zeros in \(\mathcal U_{\mathbb C}\), that the logarithm
branches are chosen compatibly, and that \(p_\Lambda\) converges uniformly on
compact subsets of \(\mathcal U_{\mathbb C}\).  Then the limiting pressure
\(p\) is holomorphic on \(\mathcal U_{\mathbb C}\), hence real analytic near
\(x_0\).  Therefore a pressure nonanalyticity at \(x_0\) requires either
Lee--Yang zeros approaching \(x_0\) in every complex neighborhood, or failure
of the stated uniform analytic convergence.
\end{proposition}

\begin{proof}
By Proposition~\ref{prop:qcd-finite-regulator-analyticity}, each
\(p_\Lambda\) is holomorphic on \(\mathcal U_{\mathbb C}\) after the stated
choice of logarithm.  The uniform compact convergence hypothesis then makes
the limit holomorphic by the Weierstrass theorem for holomorphic functions.
Restricting a holomorphic function to the real slice gives a real-analytic
function.  The contrapositive gives the final statement.
\end{proof}

\begin{proposition}[Source curvature as an integrated connected correlator]
\label{prop:qcd-source-curvature-susceptibility}
Let \(h\) be a real source coupled at finite regulator to an integrated
gauge-invariant local observable
\[
  X_\Lambda=\sum_{x\in\Lambda_s}a^3\,\mathcal O_x
\]
through the Euclidean weight \(\exp(hX_\Lambda)\).  With
\[
  p_\Lambda(h)=\frac{T}{V_\Lambda}\log Z_\Lambda(h),
  \qquad
  \overline{\mathcal O}_\Lambda=V_\Lambda^{-1}X_\Lambda,
\]
one has
\begin{equation}
  \frac{\partial p_\Lambda}{\partial h}
  =
  T\langle\overline{\mathcal O}_\Lambda\rangle_h,
  \qquad
  \frac{\partial^2 p_\Lambda}{\partial h^2}
  =
  T\,V_\Lambda
  \left(
    \langle\overline{\mathcal O}_\Lambda^2\rangle_h
    -
    \langle\overline{\mathcal O}_\Lambda\rangle_h^2
  \right).
  \label{eq:qcd-source-curvature-susceptibility}
\end{equation}
Thus a divergent susceptibility is a statement about the scaling of an
integrated connected correlator in the thermodynamic limit.
\end{proposition}

\begin{proof}
Differentiating the finite integral gives
\[
  \partial_h\log Z_\Lambda(h)=\langle X_\Lambda\rangle_h .
\]
A second derivative differentiates both the numerator and denominator of the
normalized expectation value:
\[
  \partial_h^2\log Z_\Lambda(h)
  =
  \langle X_\Lambda^2\rangle_h-\langle X_\Lambda\rangle_h^2 .
\]
Multiplying by \(T/V_\Lambda\) and writing
\(X_\Lambda=V_\Lambda\overline{\mathcal O}_\Lambda\) gives
\eqref{eq:qcd-source-curvature-susceptibility}.
\end{proof}

For physical QCD with fundamental quarks of nonzero mass, the most familiar
finite-temperature change near \(T\sim\Lambda_{\rm QCD}\) is not attached to
an exact fundamental center symmetry and not attached to an exact chiral
symmetry.  It is therefore a crossover unless a nonanalyticity is established
in the thermodynamic continuum theory.  Pure Yang--Mills theory and chiral
limits have sharper order parameters, and those are treated first.

\section{Center Deconfinement in Pure Yang--Mills}

Let \(G=SU(N_c)\) and first remove all dynamical fundamental matter.  The
thermal Wilson line at spatial point \(\mathbf x\) is
\begin{equation}
  \mathcal P(\mathbf x)
  =
  \mathcal P_\tau
  \exp\left(\ii\int_0^\beta A_0(\tau,\mathbf x)\,d\tau\right).
  \label{eq:qcd-phase-polyakov-holonomy}
\end{equation}
For an irreducible representation \(R\) define the normalized Polyakov loop
\begin{equation}
  P_R(\mathbf x)=\frac{1}{d_R}\operatorname{tr}_R\mathcal P(\mathbf x).
  \label{eq:qcd-phase-polyakov-loop}
\end{equation}
If \(R\) has \(N\)-ality \(k_R\), then the thermal center transformation by
\(z\in Z(SU(N_c))\) sends
\[
  P_R(\mathbf x)\longmapsto z^{k_R}P_R(\mathbf x).
\]
The lattice construction of the transformation and the static-source
interpretation were given in
Chapter~\ref{ch:lattice-yang-mills-nonperturbative-definition}; here we use
them as the phase-ordering language of the thermal theory.

\begin{definition}[Center-order deconfinement]
\label{def:qcd-center-order-deconfinement}
In pure \(SU(N_c)\) Yang--Mills, choose a small source \(h\) that selects a
center sector and set
\[
  \overline P_\square
  =
  V^{-1}\int_{\mathbb R^3}P_\square(\mathbf x)\,d^3x
\]
in a finite box before the infinite-volume limit.  The center order
parameter is
\begin{equation}
  \ell(T)
  =
  \lim_{h\downarrow0}\lim_{V\to\infty}
  \left\langle \overline P_\square\right\rangle_{T,h,V}.
  \label{eq:qcd-phase-polyakov-order}
\end{equation}
The thermal state is center symmetric when \(\ell(T)=0\) and center broken
when \(\ell(T)\ne0\).  A center deconfinement transition is a nonanalytic
change of the infinite-volume free energy across which this symmetry
realization changes.
\end{definition}

\begin{proposition}[Center symmetry forces the finite-volume one-point loop]
\label{prop:qcd-center-finite-volume-loop-zero}
At finite spatial volume and zero center-breaking source in pure
\(SU(N_c)\) Yang--Mills,
\[
  \left\langle P_R(\mathbf x)\right\rangle=0
\]
for every representation \(R\) whose \(N\)-ality is nonzero.
\end{proposition}

\begin{proof}
The finite-volume regulated measure and action are invariant under the
thermal center transformation.  Averaging the expectation value over the
finite group \(Z(SU(N_c))\) gives
\[
  \langle P_R\rangle
  =
  \frac1{N_c}\sum_{j=0}^{N_c-1}
  e^{2\pi\ii j k_R/N_c}\langle P_R\rangle .
\]
If \(k_R\not\equiv0\pmod{N_c}\), the finite geometric sum vanishes.  Hence
\(\langle P_R\rangle=0\).  A nonzero Polyakov-loop order parameter is
therefore necessarily a spontaneous-symmetry-breaking limit, not a
finite-volume one-point function.
\end{proof}

After multiplicative line renormalization, the one-point function has the
static-source interpretation
\begin{equation}
  \langle P_R\rangle_{\rm ren}
  =
  \exp\{-\beta(F_R-F_0)\}.
  \label{eq:qcd-phase-static-source-free-energy}
\end{equation}
Thus \(\ell=0\) is equivalent to an infinite free-energy cost for an isolated
fundamental external charge, in the center-symmetric limiting state.  This
statement is logically distinct from the zero-temperature Wilson-loop area
law.  The latter diagnoses the large-distance static potential of a
source-antisource pair; the former diagnoses the realization of thermal
center symmetry and the free energy of one wrapped external line.

\begin{remark}[Dynamical fundamental quarks]
\label{rem:qcd-dynamical-quarks-center}
Fundamental quarks explicitly break the fundamental thermal center symmetry.
Consequently \(P_\square\) remains a useful renormalized static-source
observable, but not a strict order parameter.  In physical QCD the rapid
change of the Polyakov loop is evidence for liberation of color degrees of
freedom in a crossover regime; it is not, by itself, a theorem that an exact
phase transition exists.
\end{remark}

\section{Controlled High-Temperature Holonomy Potential}
\label{sec:qcd-weiss-polyakov-holonomy-potential}

The Polyakov loop becomes a calculable order-coordinate in a regime where a
constant temporal holonomy is a valid background-field variable.  Choose a
spatially constant commuting background on
\(S^1_\beta\times\mathbb R^3\) and write its fundamental holonomy as
\begin{equation}
  P=\operatorname{diag}\left(e^{\ii\theta_1},\ldots,e^{\ii\theta_{N_c}}\right),
  \qquad
  \sum_{j=1}^{N_c}\theta_j\equiv0\pmod {2\pi}.
  \label{eq:qcd-weiss-holonomy-eigenphases}
\end{equation}
The variables \(\theta_j\) are defined modulo permutations and common shifts
compatible with \(\det P=1\).  A center transformation sends
\(P\mapsto zP\), \(z^{N_c}=1\).

\begin{controlledapproximation}[One-loop holonomy regime]
\label{ca:qcd-one-loop-holonomy-regime}
The one-loop holonomy potential is the leading weak-coupling thermal
effective potential for the zero Matsubara mode of \(A_0\), computed in a
background whose spatial gradients and noncommuting components are set to
zero.  Its expansion parameter is the running gauge coupling at scale
\(T\), and its infrared completion still requires the electric and magnetic
static sectors discussed in
Section~\ref{sec:qcd-phase-hard-thermal-loops-linde} below.
\end{controlledapproximation}

\begin{theorem}[One-loop Polyakov-holonomy potential]
\label{thm:qcd-weiss-potential}
For pure \(SU(N_c)\) Yang--Mills, the one-loop gluonic free-energy density
as a function of the holonomy \(P\) is
\begin{equation}
  f_g(P)
  =
  -\frac{2}{\pi^2\beta^4}
  \sum_{n=1}^{\infty}
  \frac{1}{n^4}\operatorname{tr}_{\rm adj}(P^n)
  =
  -\frac{2}{\pi^2\beta^4}
  \sum_{n=1}^{\infty}
  \frac{|\operatorname{tr}_{\square}P^n|^2-1}{n^4}.
  \label{eq:qcd-weiss-gluon-potential}
\end{equation}
For \(N_f\) massless Dirac fundamental quarks at \(\mu_q=0\), the one-loop
fermion contribution is
\begin{equation}
  f_q(P)
  =
  \frac{2N_f}{\pi^2\beta^4}
  \sum_{n=1}^{\infty}\frac{(-1)^n}{n^4}
  \left(
    \operatorname{tr}_{\square}P^n+
    \operatorname{tr}_{\square}P^{-n}
  \right).
  \label{eq:qcd-weiss-quark-potential}
\end{equation}
\end{theorem}

\begin{proof}
Consider first one bosonic oscillator with thermal phase \(\varphi\), meaning
that its Matsubara frequencies are shifted by the holonomy eigenvalue
\(e^{\ii\varphi}\).  Its thermal free-energy density is
\[
  T\int\frac{d^3p}{(2\pi)^3}
  \log\left(1-e^{-\beta |\mathbf p|+\ii\varphi}\right).
\]
Expanding the logarithm and integrating term by term gives
\[
  -T\sum_{n\ge1}\frac{e^{\ii n\varphi}}{n}
  \int\frac{d^3p}{(2\pi)^3}e^{-n\beta|\mathbf p|}
  =
  -\frac{1}{\pi^2\beta^4}
  \sum_{n\ge1}\frac{e^{\ii n\varphi}}{n^4},
  \label{eq:qcd-weiss-boson-oscillator}
\]
because
\[
  \int\frac{d^3p}{(2\pi)^3}e^{-n\beta|\mathbf p|}
  =
  \frac{1}{2\pi^2}\int_0^\infty p^2e^{-n\beta p}\,dp
  =
  \frac{1}{\pi^2n^3\beta^3}.
  \label{eq:qcd-weiss-thermal-momentum-integral}
\]
In background gauge, ghosts and unphysical vector polarizations cancel so
that two physical bosonic polarizations remain for each adjoint weight.
Summing \eqref{eq:qcd-weiss-boson-oscillator} over adjoint weights gives the
first expression in \eqref{eq:qcd-weiss-gluon-potential}.  The identity
\[
  \operatorname{tr}_{\rm adj}(P^n)
  =
  \operatorname{tr}_{\square}(P^n)
  \operatorname{tr}_{\overline\square}(P^n)-1
  =
  |\operatorname{tr}_{\square}P^n|^2-1
\]
uses \(\square\otimes\overline\square=\mathbf1\oplus{\rm adj}\).

For a Dirac fundamental quark at zero chemical potential the two spin
polarizations and the quark-antiquark pair give
\[
  -2T\int\frac{d^3p}{(2\pi)^3}
  \left[
    \log\left(1+e^{-\beta|\mathbf p|+\ii\theta_j}\right)
    +
    \log\left(1+e^{-\beta|\mathbf p|-\ii\theta_j}\right)
  \right]
\]
for each color eigenphase \(\theta_j\).  The plus sign inside the logarithm
is the antiperiodic thermal boundary condition.  Expanding
\(\log(1+x)=\sum_{n\ge1}(-1)^{n+1}x^n/n\), using
\eqref{eq:qcd-weiss-thermal-momentum-integral}, and summing over colors and
flavors gives \eqref{eq:qcd-weiss-quark-potential}.
\end{proof}

\begin{proposition}[Perturbative center-broken minima in pure Yang--Mills]
\label{prop:qcd-weiss-center-broken-minima}
For pure \(SU(N_c)\) Yang--Mills, the one-loop potential
\eqref{eq:qcd-weiss-gluon-potential} is minimized exactly at center
holonomies
\[
  P=z\,\mathbf1_{N_c},
  \qquad
  z^{N_c}=1 .
  \label{eq:qcd-weiss-center-minima}
\]
At these minima
\[
  f_g(z\mathbf1)
  =
  -\frac{\pi^2T^4}{45}(N_c^2-1),
  \label{eq:qcd-weiss-broken-free-energy}
\]
which is the free gluon free-energy density.
\end{proposition}

\begin{proof}
For every unitary \(N_c\times N_c\) matrix,
\[
  |\operatorname{tr}_{\square}P^n|\le N_c .
\]
Hence \(\operatorname{tr}_{\rm adj}(P^n)\le N_c^2-1\) for every \(n\).  All
coefficients in \eqref{eq:qcd-weiss-gluon-potential} are negative, so the
minimum is obtained by maximizing each
\(\operatorname{tr}_{\rm adj}(P^n)\).  The \(n=1\) term already forces
\(|\operatorname{tr}_{\square}P|=N_c\), which for a unitary matrix means all
eigenvalues are equal.  The determinant-one condition then gives
\(P=z\mathbf1_{N_c}\) with \(z^{N_c}=1\).  Substitution in
\eqref{eq:qcd-weiss-gluon-potential} gives
\[
  f_g(z\mathbf1)
  =
  -\frac{2}{\pi^2\beta^4}(N_c^2-1)
  \sum_{n\ge1}\frac1{n^4}
  =
  -\frac{2}{\pi^2\beta^4}(N_c^2-1)\frac{\pi^4}{90},
\]
which is \eqref{eq:qcd-weiss-broken-free-energy}.
\end{proof}

\begin{example}[Uniform center-symmetric holonomy]
\label{ex:qcd-weiss-center-symmetric-holonomy}
Let \(P_{\rm cs}\) have eigenvalues uniformly spaced around the unit circle,
with an overall phase chosen so that \(P_{\rm cs}\in SU(N_c)\).  Then
\[
  \operatorname{tr}_{\square}P_{\rm cs}^n=0
  \quad(N_c\nmid n),
  \qquad
  |\operatorname{tr}_{\square}P_{\rm cs}^n|^2=N_c^2
  \quad(N_c\mid n).
\]
Therefore
\[
  f_g(P_{\rm cs})
  =
  \frac{\pi^2T^4}{45}\left(1-\frac1{N_c^2}\right),
  \qquad
  f_g(P_{\rm cs})-f_g(\mathbf1)
  =
  \frac{\pi^2T^4}{45}\left(N_c^2-\frac1{N_c^2}\right).
  \label{eq:qcd-weiss-center-symmetric-cost}
\]
The perturbative high-temperature gluon determinant therefore energetically
selects a center sector rather than the uniformly distributed holonomy.
\end{example}

\begin{remark}[Fundamental quarks as a center source]
\label{rem:qcd-weiss-fundamental-quark-center-source}
The gluon potential depends on \(P\) through adjoint traces and is invariant
under \(P\mapsto zP\).  The quark potential
\eqref{eq:qcd-weiss-quark-potential} contains fundamental traces:
\[
  \operatorname{tr}_{\square}(zP)^n=z^n\operatorname{tr}_{\square}P^n .
\]
Thus dynamical fundamental quarks supply an explicit center-breaking source.
At \(P=\mathbf1\), \eqref{eq:qcd-weiss-quark-potential} gives
\[
  f_q(\mathbf1)
  =
  -\frac{7\pi^2T^4}{180}N_cN_f,
  \label{eq:qcd-weiss-free-quark-free-energy}
\]
the negative of the free massless Dirac-quark pressure in
Proposition~\ref{prop:qcd-stefan-boltzmann-pressure}.
\end{remark}

\section{Chiral Symmetry and the Dirac Spectrum}

Now set \(\mu_B=0\) and take \(N_f\) Dirac quarks with a common mass \(m\).
In the chiral limit \(m=0\), and ignoring the anomalous axial \(U(1)\), the
classical flavor symmetry is
\[
  SU(N_f)_L\times SU(N_f)_R\times U(1)_B .
\]
Let \(V_3\) be the spatial volume and let
\[
  \mathcal V_4=\beta V_3
\]
be the Euclidean thermal volume.  For degenerate quark mass, the pressure
derivative is the flavor-summed scalar source response.  We use the positive
per-flavor condensate convention
\begin{equation}
  \Sigma_m(T)
  =
  \frac1{N_f}\frac{\partial p(T,m)}{\partial m}.
  \label{eq:qcd-phase-chiral-condensate-source}
\end{equation}
Equivalently, after integrating the fermions in a fixed gauge background,
\(\Sigma_m(T)\) is the thermal average of
\(\mathcal V_4^{-1}\operatorname{Tr}(D_E+m)^{-1}\) for the one-flavor
Euclidean Dirac operator \(D_E\).  With the Grassmann-ordering convention in
which the local scalar density is
\(\bar q q=\sum_{i=1}^{N_f}\bar q_iq_i\), its flavor-summed expectation is
\[
  \langle\bar q q\rangle_{\mathrm{sum},T,m}
  =
  -N_f\Sigma_m(T).
\]
Thus \(\Sigma_m\ge0\) in a phase with the usual sign of the scalar condensate.
The chiral order parameter is the limit \(\Sigma(T)=\lim_{m\downarrow0}
\lim_{V_3\to\infty}\Sigma_m(T)\), when this limit exists.  A nonzero
condensate in the chiral limit transforms as
\((\overline{\mathbf N_f},\mathbf N_f)\oplus
(\mathbf N_f,\overline{\mathbf N_f})\) and breaks
\[
  SU(N_f)_L\times SU(N_f)_R
  \longrightarrow
  SU(N_f)_V .
\]

\begin{definition}[Chiral order parameter]
\label{def:qcd-chiral-order-parameter}
The chiral condensate used as an order parameter is
\begin{equation}
  \Sigma(T)
  =
  \lim_{m\downarrow0}\lim_{V_3\to\infty}
  \frac{1}{\mathcal V_4}
  \left\langle \operatorname{Tr}(D_E+m)^{-1}\right\rangle_{V_3,T,m},
  \label{eq:qcd-phase-sigma-definition}
\end{equation}
where \(D_E\) is the one-flavor massless Euclidean Dirac operator in a gauge
background at \(\mu_B=0\).  The infinite-volume limit precedes the chiral
limit.
\end{definition}

\begin{theorem}[Banks--Casher relation under a spectral-density hypothesis]
\label{thm:qcd-banks-casher}
Assume that the positive Dirac eigenvalue density
\[
  \rho_{\mathcal V_4}(\lambda)
  =
  \frac1{\mathcal V_4}\left\langle
  \sum_{\lambda_n>0}\delta(\lambda-\lambda_n)
  \right\rangle
\]
has a thermodynamic limit \(\rho(\lambda)\) that is continuous at
\(\lambda=0\), and that zero-mode contributions have zero density in the
same limit.  Then
\begin{equation}
  \Sigma(T)=\pi\rho(0)
  \label{eq:qcd-phase-banks-casher}
\end{equation}
with the normalization in
\eqref{eq:qcd-phase-sigma-definition}.  If \(\rho(0)=0\), this particular
condensate vanishes.
\end{theorem}

\begin{proof}
At \(\mu_B=0\), \(D_E\) is anti-Hermitian in the usual positive Euclidean
inner product and anticommutes with \(\gamma_5\) in the massless theory.
Nonzero eigenvalues occur in pairs \(\pm\ii\lambda_n\), \(\lambda_n>0\).
For fixed positive mass,
\[
  \frac1{\mathcal V_4}\operatorname{Tr}(D_E+m)^{-1}
  =
  \frac{N_0}{\mathcal V_4m}
  +
  \frac1{\mathcal V_4}\sum_{\lambda_n>0}
  \left(\frac1{m+\ii\lambda_n}+\frac1{m-\ii\lambda_n}\right),
\]
where \(N_0\) is the number of zero modes.  The nonzero pair contribution is
\[
  \frac{2m}{\mathcal V_4}\sum_{\lambda_n>0}\frac1{\lambda_n^2+m^2}.
\]
By the zero-density hypothesis the first term drops out after
\(V_3\to\infty\).  The condensate definition therefore gives
\[
  \Sigma(T)
  =
  \lim_{m\downarrow0}
  \int_0^\infty
  \frac{2m\,\rho(\lambda)}{\lambda^2+m^2}\,d\lambda .
\]
The kernels
\[
  K_m(\lambda)=\frac{2m}{\lambda^2+m^2},\qquad \lambda\ge0,
\]
form an approximate identity on the half-line with total mass \(\pi\):
\[
  \int_0^\infty K_m(\lambda)\,d\lambda=\pi .
\]
Continuity of \(\rho\) at the origin gives the limit
\(\pi\rho(0)\), and hence \eqref{eq:qcd-phase-banks-casher}.  If instead
one defines the spectral density on all of \(\mathbb R\), counting both
positive and negative eigenvalues, the same formula holds with that
two-sided density evaluated at zero.
\end{proof}

\begin{remark}[Order of limits]
The theorem is a statement about the thermodynamic spectral density near
zero.  Setting \(m=0\) at finite volume first gives no condensate from paired
nonzero eigenvalues.  The density of eigenvalues at spacing
\(O(\mathcal V_4^{-1})\) near the origin is precisely the infinite-volume
datum that can survive the subsequent chiral limit.
\end{remark}

\subsection{Axial Ward Identity and Thermal GMOR Logic}
\label{subsec:qcd-thermal-axial-ward-gmor}

The condensate is not merely a symbol for a phase diagram.  It is tied to
integrated pseudoscalar correlations by an axial Ward identity.  This is the
cleanest finite-temperature form of the Gell-Mann--Oakes--Renner logic,
because at \(T>0\) Lorentz invariance is absent and one must specify which
pion pole, screening pole, or spectral residue is being used.

Let \(\tau^a\), \(a=1,\ldots,N_f^2-1\), be Hermitian generators of
\(\mathfrak{su}(N_f)\) normalized by
\begin{equation}
  \operatorname{tr}_{f}(\tau^a\tau^b)=\delta^{ab}.
  \label{eq:qcd-flavor-generator-normalization}
\end{equation}
Define the nonsinglet Euclidean axial current and pseudoscalar density by
\begin{equation}
  A_\mu^a=\bar q\gamma_\mu\gamma_5\tau^a q,
  \qquad
  P^a=\bar q\,\ii\gamma_5\tau^a q .
  \label{eq:qcd-axial-current-pseudoscalar-density}
\end{equation}
The factor of \(\ii\) in \(P^a\) is part of the Euclidean source convention:
a real pseudoscalar source \(p^a\) appears in the mass matrix as
\(m\mathbf1+\ii\gamma_5p^a\tau^a\).  With this convention the Euclidean
pseudoscalar susceptibility is
\begin{equation}
  \chi_\pi^{ab}(T,m)
  =
  -\int_0^\beta d\tau\int_{\mathbb R^3}d^3x\,
  \langle P^a(\tau,\mathbf x)P^b(0)\rangle_{c,T,m},
  \label{eq:qcd-pseudoscalar-susceptibility-definition}
\end{equation}
with the finite-volume definition obtained by replacing \(\mathbb R^3\) by
the spatial torus and taking the thermodynamic limit after the integral.

\begin{proposition}[Integrated nonsinglet axial Ward identity]
\label{prop:qcd-integrated-axial-ward}
Assume a regulator and renormalization prescription preserving the nonsinglet
vector flavor symmetry and the nonsinglet axial Ward identity up to the
explicit mass term; there is no nonsinglet axial anomaly.  Then
\begin{equation}
  2m\,\chi_\pi^{ab}(T,m)
  =
  -\left\langle
  \bar q\{\tau^a,\tau^b\}q
  \right\rangle_{T,m}.
  \label{eq:qcd-integrated-axial-ward}
\end{equation}
In a flavor-symmetric thermal state this becomes
\begin{equation}
  m\,\chi_\pi^{ab}(T,m)=\delta^{ab}\Sigma_m(T).
  \label{eq:qcd-flavor-symmetric-pseudoscalar-ward}
\end{equation}
\end{proposition}

\begin{proof}
Perform the local nonsinglet axial change of variables
\[
  \delta q=\ii\alpha^a(x)\tau^a\gamma_5 q,\qquad
  \delta\bar q=\ii\alpha^a(x)\bar q\gamma_5\tau^a .
\]
The regulator hypothesis says that the fermion measure has no nonsinglet
Jacobian anomaly.  The action variation is the distributional Ward
insertion
\[
  \partial_\mu A_\mu^a(x)-2mP^a(x),
\]
with contact terms when \(x\) collides with other insertions.  Apply this
identity to the insertion \(P^b(0)\).  The contact term is its infinitesimal
axial variation:
\[
  \delta^a P^b(0)
  =
  -\bar q\{\tau^a,\tau^b\}q(0).
\]
Integrating the Ward identity over the thermal spacetime torus removes the
divergence term by periodicity of the bosonic current and gives
\eqref{eq:qcd-integrated-axial-ward}, with the minus sign in
\eqref{eq:qcd-pseudoscalar-susceptibility-definition} coming from the
Euclidean \(P^a=\bar q\ii\gamma_5\tau^a q\) convention.

In a flavor-symmetric state
\[
  \left\langle\bar q\{\tau^a,\tau^b\}q\right\rangle_{T,m}
  =
  \frac{2}{N_f}\delta^{ab}
  \langle\bar q q\rangle_{\mathrm{sum},T,m}
  =
  -2\delta^{ab}\Sigma_m(T),
\]
where \eqref{eq:qcd-flavor-generator-normalization} fixes the singlet part of
the anticommutator.  Substitution gives
\eqref{eq:qcd-flavor-symmetric-pseudoscalar-ward}.
\end{proof}

\begin{proposition}[GMOR as a pole-saturation consequence]
\label{prop:qcd-thermal-gmor-pole-saturation}
Suppose, in addition to the hypotheses of
Proposition~\ref{prop:qcd-integrated-axial-ward}, that for small \(m>0\) the
zero-spatial-momentum pseudoscalar susceptibility in a broken phase is
saturated at leading order by an isolated pion pole of mass \(M_\pi(T,m)\).
Let \(G_\pi(T,m)\) and \(f_{t,\delta}(T,m)\) be defined in the same flavor
normalization \eqref{eq:qcd-flavor-generator-normalization} by the pole
matrix elements
\[
  \langle\Omega_T|P^a(0)|\pi^b\rangle
  =
  \delta^{ab}G_\pi,
  \qquad
  \langle\Omega_T|A_0^a(0)|\pi^b\rangle
  =
  \delta^{ab}f_{t,\delta}M_\pi,
\]
or by the corresponding KMS spectral residues when a thermal Hilbert-space
representative is used.  If the pole obeys the axial Ward relation
\begin{equation}
  2m\,G_\pi=f_{t,\delta}M_\pi^2
  \label{eq:qcd-pcac-pole-relation}
\end{equation}
and the non-pole part of \(\chi_\pi^{aa}\) is \(O(1)\) as
\(m\downarrow0\), then
\begin{equation}
  f_{t,\delta}(T,m)^2M_\pi(T,m)^2
  =
  4m\Sigma(T)+o(m).
  \label{eq:qcd-thermal-gmor-delta-normalization}
\end{equation}
If instead one uses generators \(t^a\) normalized by
\(\operatorname{tr}_{f}(t^at^b)=\frac12\delta^{ab}\), the decay constant is
smaller by \(\sqrt2\) and the same statement is
\begin{equation}
  f_{t,1/2}(T,m)^2M_\pi(T,m)^2
  =
  2m\Sigma(T)+o(m).
  \label{eq:qcd-thermal-gmor-half-normalization}
\end{equation}
\end{proposition}

\begin{proof}
Pole saturation gives
\[
  \chi_\pi^{ab}(T,m)
  =
  \delta^{ab}\frac{G_\pi(T,m)^2}{M_\pi(T,m)^2}+O(1).
\]
The Ward identity \eqref{eq:qcd-flavor-symmetric-pseudoscalar-ward} says that
the same susceptibility has leading behavior
\[
  \chi_\pi^{ab}(T,m)
  =
  \delta^{ab}\frac{\Sigma(T)}{m}+o(m^{-1})
\]
when \(\Sigma_m(T)\to\Sigma(T)\).  Using
\eqref{eq:qcd-pcac-pole-relation} to replace \(G_\pi\) gives
\[
  \frac{G_\pi^2}{M_\pi^2}
  =
  \frac{f_{t,\delta}^2M_\pi^2}{4m^2}.
\]
Equating the coefficients of \(1/m\) yields
\eqref{eq:qcd-thermal-gmor-delta-normalization}.  Rescaling the generators
from \(\operatorname{tr}_f(\tau^a\tau^b)=\delta^{ab}\) to
\(\operatorname{tr}_f(t^at^b)=\frac12\delta^{ab}\) rescales both the current
and the pseudoscalar density by \(1/\sqrt2\), hence
\(f_{t,1/2}=f_{t,\delta}/\sqrt2\), giving
\eqref{eq:qcd-thermal-gmor-half-normalization}.
\end{proof}

\subsection{Low-Temperature Chiral Effective Theory}
\label{subsec:qcd-low-temperature-chiral-effective-theory}

The Ward identity identifies the source derivative of the pressure with the
integrated pion channel.  In the low-temperature phase this identity can be
evaluated in a controlled effective theory.  The control assumptions are
part of the statement: the spatial thermodynamic limit has been taken, the
temperature is small compared with the first non-Goldstone mass scale, and
the derivative expansion is organized at momenta \(p\ll 4\pi F\).

\begin{definition}[Leading chiral effective datum]
\label{def:qcd-leading-chiral-effective-datum}
For \(N_f\) degenerate light quarks, the leading Euclidean chiral effective
datum consists of a field \(U:\mathbb R^3\times S^1_\beta\to SU(N_f)\), a
constant \(F>0\), and a constant \(B>0\), with leading action
\begin{equation}
  S_{\chi}^{(2)}
  =
  \int_0^\beta d\tau\int d^3x
  \left[
    \frac{F^2}{4}
    \operatorname{tr}_f(\partial_\mu U^\dagger\partial_\mu U)
    -
    \frac{F^2B}{2}
    \operatorname{tr}_f\bigl(MU^\dagger+UM^\dagger\bigr)
  \right].
  \label{eq:qcd-leading-chiral-effective-action}
\end{equation}
Here \(M=m\mathbf1_{N_f}\) for the degenerate-mass discussion, and
\(\operatorname{tr}_f\) is the ordinary fundamental flavor trace.  We
parameterize the pion fields by
\begin{equation}
  U=\exp\left(\frac{\ii\sqrt2}{F}\pi^a\tau^a\right),
  \qquad
  \operatorname{tr}_f(\tau^a\tau^b)=\delta^{ab}.
  \label{eq:qcd-chiral-field-parameterization}
\end{equation}
With this convention \(F\) is the decay constant in the half-trace flavor
normalization; the trace-delta decay constant used in
Proposition~\ref{prop:qcd-thermal-gmor-pole-saturation} is
\(f_{t,\delta}=\sqrt2F+O(m,T^2)\).
\end{definition}

\begin{proposition}[Tree-level source normalization]
\label{prop:qcd-chiral-effective-source-normalization}
The datum \eqref{eq:qcd-leading-chiral-effective-action} gives
\begin{equation}
  \Sigma=F^2B,
  \qquad
  M_\pi^2=2Bm
  \label{eq:qcd-chiral-sigma-mpi-tree}
\end{equation}
at tree level, where \(\Sigma\) is the positive per-flavor condensate in
\eqref{eq:qcd-phase-chiral-condensate-source}.  Equivalently,
\[
  F^2M_\pi^2=2m\Sigma,
  \qquad
  f_{t,\delta}^2M_\pi^2=4m\Sigma
\]
at this order.
\end{proposition}

\begin{proof}
For a constant field \(U=\mathbf1\), the Euclidean free-energy density from
\eqref{eq:qcd-leading-chiral-effective-action} is
\[
  f_0(m)=-N_fF^2Bm,
  \qquad
  p_0(m)=N_fF^2Bm .
\]
Therefore \(N_f^{-1}\partial p_0/\partial m=F^2B\), proving
\(\Sigma=F^2B\).  Expanding
\eqref{eq:qcd-chiral-field-parameterization},
\[
  U+U^\dagger
  =
  2\mathbf1-\frac{2}{F^2}\pi^a\pi^b\tau^a\tau^b+O(\pi^4),
\]
and substituting \(M=m\mathbf1\) gives the quadratic Euclidean Lagrangian
\[
  \frac12\partial_\mu\pi^a\partial_\mu\pi^a
  +
  Bm\,\pi^a\pi^a .
\]
Comparing with
\(\frac12(\partial\pi)^2+\frac12M_\pi^2\pi^2\) gives
\(M_\pi^2=2Bm\).  The two GMOR forms follow from
\(f_{t,\delta}=\sqrt2F\) at leading order.
\end{proof}

\begin{proposition}[One-loop pion-gas correction to the condensate]
\label{prop:qcd-low-temperature-chiral-condensate}
Under the leading chiral effective assumptions above, and in the chiral limit
after the spatial thermodynamic limit, the low-temperature condensate obeys
\begin{equation}
  \frac{\Sigma(T)}{\Sigma(0)}
  =
  1-\frac{N_f^2-1}{12N_f}\frac{T^2}{F^2}
  +O\!\left(\frac{T^4}{F^4}\right)
  \label{eq:qcd-low-temperature-chiral-condensate}
\end{equation}
within the chiral derivative expansion.  For \(N_f=2\), this is
\[
  \Sigma(T)/\Sigma(0)=1-\frac{T^2}{8F^2}+O(T^4/F^4).
\]
\end{proposition}

\begin{proof}
To one loop in the leading effective theory the only \(m\)-dependent thermal
term at order \(T^2\) is the free pressure of
\(d_\pi=N_f^2-1\) pion fields with
\(\omega_k=(k^2+M_\pi^2)^{1/2}\):
\begin{equation}
  p_\pi(T,M_\pi)
  =
  -d_\pi T\int\frac{d^3k}{(2\pi)^3}
  \log\bigl(1-e^{-\beta\omega_k}\bigr).
  \label{eq:qcd-free-pion-pressure}
\end{equation}
The vacuum part is absorbed into the zero-temperature low-energy constants;
the displayed thermal part is finite.  Differentiating at fixed \(T\),
\begin{equation}
  \frac{\partial p_\pi}{\partial M_\pi^2}
  =
  -\frac{d_\pi}{2}
  \int\frac{d^3k}{(2\pi)^3}
  \frac{1}{\omega_k}\frac{1}{e^{\beta\omega_k}-1}.
  \label{eq:qcd-pion-pressure-mass-derivative}
\end{equation}
Since \(M_\pi^2=2Bm\), the pion contribution to the per-flavor condensate is
\[
  \frac1{N_f}\frac{\partial p_\pi}{\partial m}
  =
  -\frac{d_\pi B}{N_f}
  \int\frac{d^3k}{(2\pi)^3}
  \frac{1}{\omega_k}\frac{1}{e^{\beta\omega_k}-1}.
\]
In the chiral limit the integral is
\[
  \int\frac{d^3k}{(2\pi)^3}
  \frac{1}{|k|}\frac{1}{e^{\beta |k|}-1}
  =
  \frac{T^2}{12}.
  \label{eq:qcd-massless-pion-tadpole}
\]
Thus
\[
  \Sigma(T)
  =
  F^2B-\frac{N_f^2-1}{N_f}B\,\frac{T^2}{12}
  +O(T^4/F^2).
\]
Dividing by \(\Sigma(0)=F^2B\) gives
\eqref{eq:qcd-low-temperature-chiral-condensate}.  The next terms require
the \(O(p^4)\) chiral Lagrangian and two-loop pion interactions; they are
not fixed by symmetry and \(F,B\) alone.
\end{proof}

\begin{remark}[What the low-temperature formula establishes]
Equation~\eqref{eq:qcd-low-temperature-chiral-condensate} is a controlled
asymptotic statement inside the broken phase.  It does not by itself prove a
chiral restoration temperature or determine the order of a transition.  Such
questions require the nonperturbative thermal QCD limit, or an effective
description whose domain includes the critical region.
\end{remark}

\section{Free Plasma Limit and Thermodynamic Observables}

At asymptotically high temperature and fixed \(\mu_B/T\), asymptotic freedom
makes the leading term in the pressure the free massless quark-gluon pressure.
This is not a statement that the plasma is free at any finite temperature;
it is the leading term in a controlled high-temperature expansion whose
subleading terms require screening and infrared resummation.

\begin{proposition}[Massless free QCD pressure]
\label{prop:qcd-stefan-boltzmann-pressure}
For \(SU(N_c)\) with \(N_f\) massless Dirac fundamental quarks at
\(\mu_B=0\), the free pressure is
\begin{equation}
  p_{\rm SB}(T)
  =
  \frac{\pi^2T^4}{45}(N_c^2-1)
  +
  \frac{7\pi^2T^4}{180}N_cN_f .
  \label{eq:qcd-phase-stefan-boltzmann}
\end{equation}
At nonzero baryon chemical potential, the quark part becomes
\begin{equation}
  p_{q,{\rm free}}(T,\mu_B)
  =
  N_cN_f\left[
  \frac{7\pi^2T^4}{180}
  +\frac{\mu_B^2T^2}{6N_c^2}
  +\frac{\mu_B^4}{12\pi^2N_c^4}
  \right].
  \label{eq:qcd-phase-free-quark-mu-pressure}
\end{equation}
\end{proposition}

\begin{proof}
A massless bosonic degree of freedom contributes
\[
  -T\int\frac{d^3p}{(2\pi)^3}\log(1-e^{-|p|/T})
  =
  \frac{\pi^2T^4}{90},
\]
using \(\int_0^\infty x^3(e^x-1)^{-1}dx=\pi^4/15\).  A gluon has two
physical polarizations and \(N_c^2-1\) color states, giving the first term.

A massless fermionic degree of freedom contributes \(7/8\) times the bosonic
answer at zero chemical potential.  A Dirac quark has four
particle-antiparticle and spin degrees for each color and flavor, giving
\((7/8)(4N_cN_f)\pi^2T^4/90\), which is the second term of
\eqref{eq:qcd-phase-stefan-boltzmann}.  The standard Fermi integral with
chemical potential \(\mu_q\) gives, for one Dirac quark of one color,
\[
  \frac{7\pi^2T^4}{180}
  +\frac{\mu_q^2T^2}{6}
  +\frac{\mu_q^4}{12\pi^2}.
\]
Substituting \(\mu_q=\mu_B/N_c\) and multiplying by \(N_cN_f\) gives
\eqref{eq:qcd-phase-free-quark-mu-pressure}.
\end{proof}

The energy density, entropy density, and trace anomaly are then defined by
\[
  s=\frac{\partial p}{\partial T},\qquad
  n_B=\frac{\partial p}{\partial\mu_B},\qquad
  \epsilon=-p+Ts+\mu_B n_B,
\]
and
\[
  \Delta(T,\mu_B)=\epsilon-3p .
\]
For a conformal thermal state in four dimensions, \(\Delta=0\).  In QCD,
\(\Delta\) measures the breaking of scale invariance by quark masses, by the
running coupling, and by the thermal state itself.

\begin{proposition}[Thermodynamic derivative identities for QCD observables]
\label{prop:qcd-thermodynamic-derivative-identities}
Let \(p(T,\mu_B)\) be the thermodynamic pressure of a QCD phase on an open
domain in which it is twice continuously differentiable.  Define
\[
  s=\partial_Tp,\qquad n_B=\partial_{\mu_B}p,\qquad
  \epsilon=-p+Ts+\mu_Bn_B .
\]
Then
\[
  dp=s\,dT+n_B\,d\mu_B,\qquad
  d\epsilon=T\,ds+\mu_B\,dn_B,
\]
and the interaction measure is
\[
  \Delta(T,\mu_B)=\epsilon-3p
  =
  T\partial_Tp+\mu_B\partial_{\mu_B}p-4p .
\]
Equivalently, at fixed \(x=\mu_B/T\),
\[
  \frac{\Delta(T,\mu_B)}{T^4}
  =
  T\,\partial_T\left(\frac{p(T,xT)}{T^4}\right)_{x}.
\]
At \(\mu_B=0\), if charge conjugation gives \(n_B(T,0)=0\) and
\(\partial_T\epsilon(T,0)>0\), the adiabatic sound speed determined by the
pressure is
\[
  c_s^2(T)
  =
  \frac{dp(T,0)/dT}{d\epsilon(T,0)/dT}
  =
  \frac{\partial_Tp(T,0)}{T\,\partial_T^2p(T,0)} .
\]
These formulae are definitions and thermodynamic consequences of the
pressure; they are not model assumptions about quasiparticles.
\end{proposition}

\begin{proof}
The first identity is the definition of \(s\) and \(n_B\) as derivatives of
the limiting pressure.  Differentiating
\(\epsilon=-p+Ts+\mu_Bn_B\) gives
\[
  d\epsilon
  =
  -dp+T\,ds+s\,dT+\mu_B\,dn_B+n_B\,d\mu_B.
\]
Substituting \(dp=s\,dT+n_B\,d\mu_B\) leaves
\(d\epsilon=T\,ds+\mu_B\,dn_B\).  The displayed formula for \(\Delta\)
follows by substituting \(s=\partial_Tp\) and
\(n_B=\partial_{\mu_B}p\) into \(\epsilon-3p\).

For the fixed-\(x\) identity, write \(\mu_B=xT\).  The chain rule gives
\[
  T\,\partial_T\left(\frac{p(T,xT)}{T^4}\right)_{x}
  =
  \frac{T\partial_Tp+\mu_B\partial_{\mu_B}p-4p}{T^4}.
\]
At \(\mu_B=0\), charge conjugation removes the baryon-density term.  Along
that one-parameter equilibrium curve,
\[
  \frac{d\epsilon}{dT}
  =
  T\,\frac{ds}{dT}
  =
  T\,\partial_T^2p(T,0),
  \qquad
  \frac{dp}{dT}=\partial_Tp(T,0),
\]
which gives the sound-speed formula whenever the denominator is nonzero.
\end{proof}

\begin{example}[Conformal Stefan--Boltzmann benchmark]
\label{ex:qcd-conformal-sb-thermodynamics}
If
\[
  p(T,\mu_B)=aT^4+b\mu_B^2T^2+c\mu_B^4
\]
with \(a,b,c\) independent of \(T\) and \(\mu_B\), then \(p\) is homogeneous
of degree four, so Euler's identity gives
\[
  T\partial_Tp+\mu_B\partial_{\mu_B}p=4p,\qquad \Delta=0.
\]
At \(\mu_B=0\), one obtains \(\epsilon=3aT^4\) and
\[
  c_s^2=\frac{4aT^3}{12aT^3}=\frac13 .
\]
The free plasma is therefore a normalization benchmark for lattice and
resummed calculations of \(\Delta/T^4\) and \(c_s^2\); departures from it
measure interactions, masses, running couplings, and nonperturbative thermal
physics.
\end{example}

\section{Hard Thermal Loops and the Linde Obstruction}
\label{sec:qcd-phase-hard-thermal-loops-linde}

Chapter~\ref{ch:thermal-gauge-theory-screening} derived the one-loop Debye
mass in the present trace convention:
\[
  m_D^2
  =
  g^2T^2\left(
  \frac{C_A}{3}
  +\frac{1}{3}\sum_{\text{Dirac }f}T_{R_f}
  \right)
\]
for QCD without scalars.  For \(SU(N_c)\) with
\(\operatorname{tr}_{\square}(T^aT^b)=\delta^{ab}\), this gives
\[
  C_A=2N_c,\qquad T_{\square}=1.
\]
The scale \(gT\) is electric.  It must be resummed in static longitudinal
propagators before a weak-coupling expansion of many thermal quantities is
well-defined.

\begin{controlledapproximation}[Hard-thermal-loop regime]
\label{ca:qcd-htl-regime}
The HTL approximation applies to external frequencies and momenta
\[
  \omega,\ |\mathbf k| = O(gT)
\]
in a plasma whose internal loop momenta are hard, \(p=O(T)\).  The leading
soft gauge-field effective action is the gauge-invariant functional whose
second variation reproduces the transverse and longitudinal one-loop
polarization tensors in this kinematic regime.  A calculation outside this
scale separation is not an HTL calculation.
\end{controlledapproximation}

\begin{proposition}[Static HTL polarization and Debye mass]
\label{prop:qcd-static-htl-debye-mass}
Let \(a_0=a_0^aT^a\) be a slowly varying canonically normalized temporal
gauge potential, and normalize the generators by
\[
  \operatorname{tr}_R(T_R^aT_R^b)=T_R\,\delta^{ab},\qquad
  f^{acd}f^{bcd}=C_A\,\delta^{ab}.
\]
For a weakly coupled thermal plasma at \(\mu_B=0\), the static HTL color
susceptibility is
\[
  \Pi_{00}^{ab}(0,\mathbf k\to0)
  =
  \delta^{ab}m_D^2,
  \qquad
  m_D^2
  =
  g^2T^2\left(\frac{C_A}{3}
  +\frac{1}{3}\sum_{\text{Dirac }f}T_{R_f}\right).
\]
In the present \(SU(N_c)\) convention
\(\operatorname{tr}_{\square}(T^aT^b)=\delta^{ab}\), one has
\[
  C_A=2N_c,\qquad T_{\square}=1,
\]
and hence
\[
  m_D^2=g^2T^2\left(\frac{2N_c}{3}+\frac{N_f}{3}\right).
\]
To leading static order the longitudinal electric propagator is screened as
\[
  D_{00}^{ab}(0,\mathbf k)
  =
  \frac{\delta^{ab}}{\mathbf k^2+m_D^2},
\]
with the understanding that this formula is the static electric HTL limit,
not a gauge-invariant definition of a particle pole.
\end{proposition}

\begin{proof}
The statement is a color-susceptibility calculation for hard particles.
Let a hard excitation of momentum \(p=|\mathbf p|\) and color weight \(w\)
couple to the soft potential by the energy shift
\(-g a_0^aw^a\).  Linearizing its distribution gives
\[
  \delta n(p)=g a_0^bw^b\,[-n'(p)] .
\]
Multiplying by the color charge \(gw^a\), summing over weights, and
integrating over hard momenta gives the induced color density.  Therefore a
species in representation \(R\) contributes the positive susceptibility
\[
  g^2\,d_{\mathrm{spin/particle}}\,
  \operatorname{tr}_R(T_R^aT_R^b)
  \int \frac{d^3p}{(2\pi)^3}\,[-n'(p)] .
\]
For gluons the representation is the adjoint and there are two transverse
hard polarizations.  For a massless Dirac fermion there are two spin states
and independent particle and antiparticle excitations; the antiparticle
charge is opposite, but the susceptibility is quadratic in the charge, so
the total degeneracy factor is \(4\).

The required thermal integrals are elementary but convention-sensitive:
\[
  I_B
  =
  \int \frac{d^3p}{(2\pi)^3}\,[-n_B'(p)]
  =
  \frac{1}{\pi^2}\int_0^\infty dp\,p\,n_B(p)
  =\frac{T^2}{6},
\]
\[
  I_F
  =
  \int \frac{d^3p}{(2\pi)^3}\,[-n_F'(p)]
  =
  \frac{1}{\pi^2}\int_0^\infty dp\,p\,n_F(p)
  =\frac{T^2}{12}.
\]
The integration by parts has no boundary term: at \(p=\infty\) the
distributions decay exponentially, while at \(p=0\) one has
\(p^2n_B(p)\to0\) and \(p^2n_F(p)\to0\); moreover
\(p^2[-n'_B(p)]\) and \(p^2[-n'_F(p)]\) are locally integrable.  The last
equalities use
\[
  \int_0^\infty \frac{x\,dx}{e^x-1}=\frac{\pi^2}{6},
  \qquad
  \int_0^\infty \frac{x\,dx}{e^x+1}=\frac{\pi^2}{12}.
\]
Thus
\[
  m_D^2
  =
  g^2\bigl(2C_A I_B+4\sum_fT_{R_f}I_F\bigr)
  =
  g^2T^2\left(\frac{C_A}{3}
  +\frac{1}{3}\sum_fT_{R_f}\right).
\]
The resummed static longitudinal propagator is the geometric series obtained
by inserting this static self-energy into the soft electric propagator.  The
derivation has used only the HTL separation between hard loop momentum and
soft external field; it does not determine the magnetic sector.
\end{proof}

\begin{proposition}[Magnetic contribution first becomes nonperturbative at \(g^6T^4\)]
\label{prop:qcd-linde-scale}
In four-dimensional high-temperature Yang--Mills theory, the static magnetic
sector has three-dimensional coupling
\[
  g_3^2=g^2T .
\]
If the magnetic sector develops a mass gap of order \(g_3^2\), its
contribution to the four-dimensional pressure begins at order
\[
  T(g_3^2)^3 = g^6T^4 .
\]
The coefficient of this contribution is not determined by ordinary
perturbation theory around the massless magnetic theory.
\end{proposition}

\begin{proof}
After the nonzero Matsubara modes and the electric scale have been matched
into a three-dimensional static effective theory, the remaining magnetic
Yang--Mills action has the form
\[
  \frac1{4g_3^2}\int d^3x\,\operatorname{tr}(F_{ij}F_{ij}).
\]
In three dimensions \(g_3^2\) has mass dimension one.  A vacuum free-energy
density of the magnetic theory has mass dimension three.  If no scale other
than \(g_3^2\) remains in the confining magnetic sector, dimensional analysis
forces the nonperturbative contribution to be \(c(g_3^2)^3\), where \(c\) is
a dimensionless number defined by the three-dimensional Yang--Mills theory.
The four-dimensional pressure is \(T\) times the static three-dimensional
free-energy density, so the order is \(g^6T^4\).

Perturbation theory cannot determine \(c\), because it expands around a
massless three-dimensional gauge theory whose infrared behavior is precisely
the magnetic theory being asked to generate the scale.  The proposition does
not prove the magnetic mass gap; it identifies the order at which the
nonperturbative magnetic datum first enters the high-temperature expansion.
\end{proof}

\section{Transport and Quark--Gluon Plasma Observables}

The quark--gluon plasma is not defined by the presence of free quarks and
gluons.  In this monograph it is the thermal regime of QCD in which the
dominant excitations and response functions are better organized by color
degrees of freedom, screening, and hydrodynamic transport than by a dilute
gas of hadrons.  The precise observables are correlation functions.

The shear viscosity, bulk viscosity, and electrical conductivity are defined
by the Kubo limits of Chapter~\ref{ch:thermal-spectral-kubo-transport}.  For
example, with \(T_{xy}\) a spatial off-diagonal stress tensor component,
\[
  \eta
  =
  \lim_{\omega\downarrow0}
  \frac{1}{\omega}\,
  \rho_{T_{xy}T_{xy}}(\omega,\mathbf0)
\]
when the limit exists after the thermodynamic limit.  The definition is
nonperturbative; a kinetic or holographic computation is an approximation to
this spectral limit, not the definition of the coefficient.

\begin{remark}[The \(\eta/s\) lower-bound claim]
The value \(1/(4\pi)\) is an important result in a class of strongly coupled
large-\(N\) theories with classical gravity duals.  It is not a theorem of
QCD.  When this number is used as a comparison scale in QCD phenomenology,
the statement is a quantitative comparison of transport data or models, not
a rigorous lower bound derived from the QCD axioms.
\end{remark}

\section{Finite Baryon Density and the Sign Problem}

At real \(\mu_B\), the Euclidean Dirac operator is no longer anti-Hermitian
up to the mass term.  In continuum notation the temporal derivative is
shifted by \(-\mu_q\), and on the lattice forward and backward temporal hops
are weighted asymmetrically.  The determinant need not be real and
nonnegative.

\begin{proposition}[Origin of the finite-density sign problem]
\label{prop:qcd-finite-density-sign-problem}
At \(\mu_B=0\), a vectorlike pair of equal-mass quark flavors gives a
nonnegative Euclidean determinant.  At real \(\mu_B\ne0\), the same
positivity argument fails because the Dirac operator loses the
\(\gamma_5\)-Hermiticity relation that pairs eigenvalues into complex
conjugates at fixed gauge field.
\end{proposition}

\begin{proof}
At \(\mu_B=0\), the lattice or continuum Euclidean Dirac operator obeys
\[
  \gamma_5(D_E+m)\gamma_5=(D_E+m)^\dagger
\]
for real mass \(m\).  Therefore
\(\det(D_E+m)\) is real, and for two degenerate flavors the product is
\(|\det(D_E+m)|^2\ge0\).  This is the measure-positivity input behind
importance sampling.

A real quark chemical potential changes the temporal part of the operator by
a term that is Hermitian rather than anti-Hermitian in Euclidean signature.
Equivalently, the lattice temporal hopping factors \(e^{a\mu_q}\) and
\(e^{-a\mu_q}\) are not related by the same unitary conjugation that holds at
\(\mu_q=0\).  The displayed \(\gamma_5\)-Hermiticity relation is replaced by
a relation connecting \(\mu_q\) to \(-\mu_q\).  At fixed real \(\mu_q\) the
determinant is generally complex, so it cannot be used directly as a
positive probability weight.
\end{proof}

Consequently, the conjectural critical endpoint and the first-order line in
the \((T,\mu_B)\) plane are not presently consequences of direct
first-principles continuum control at large real \(\mu_B\).  Taylor
expansions around \(\mu_B=0\), imaginary-\(\mu_B\) continuation, reweighting,
and effective models are valuable but must carry their own domains of
validity.

\section{Baryon Susceptibilities and Critical-Point Diagnostics}
\label{sec:qcd-baryon-susceptibilities-criticality}

At \(\mu_B=0\) the Euclidean measure of vectorlike QCD is positive in the
standard lattice discretizations, so derivatives with respect to
\(\mu_B\) give first-principles thermal observables even though direct
importance sampling at real \(\mu_B\ne0\) is obstructed.  The natural
dimensionless source is
\[
  x=\frac{\mu_B}{T}=\beta\mu_B .
\]
At finite regulator and finite volume,
\begin{equation}
  Z_{a,V}(x)
  =
  \operatorname{Tr}_{\Hilb_{a,V}}
  \exp\{-\beta H_{a,V}+xB_{a,V}\}.
  \label{eq:qcd-baryon-finite-x-partition}
\end{equation}
Here \(B_{a,V}\) is the finite-regulator baryon charge, normalized so that a
baryon has charge \(1\), and it is assumed to be an exactly conserved
self-adjoint charge of the regulated theory:
\[
  [H_{a,V},B_{a,V}]=0 .
\]

\begin{definition}[Baryon susceptibilities]
\label{def:qcd-baryon-susceptibilities}
The finite-regulator baryon cumulants and susceptibilities at \(x=0\) are
\begin{equation}
  \kappa_{n,a,V}^{B}(T)
  =
  \left.
  \frac{\partial^n}{\partial x^n}\log Z_{a,V}(x)
  \right|_{x=0},
  \qquad
  \chi_{n,a,V}^{B}(T)
  =
  \frac{1}{VT^3}\kappa_{n,a,V}^{B}(T).
  \label{eq:qcd-baryon-susceptibility-definition}
\end{equation}
The continuum thermodynamic susceptibility \(\chi_n^B(T)\) is the limit of
\(\chi_{n,a,V}^{B}(T)\) along the chosen continuum trajectory, when the
limit exists.
\end{definition}

\begin{proposition}[Baryon susceptibilities are baryon-number cumulants]
\label{prop:qcd-baryon-cumulant-identities}
At finite regulator and finite volume,
\begin{equation}
  \chi_{1,a,V}^B
  =
  \frac{\langle B\rangle}{VT^3},
  \qquad
  \chi_{2,a,V}^B
  =
  \frac{\langle B^2\rangle-\langle B\rangle^2}{VT^3},
  \label{eq:qcd-baryon-first-second-cumulants}
\end{equation}
and
\begin{align}
  \chi_{4,a,V}^B
  &=
  \frac{1}{VT^3}
  \bigl(
  \langle B^4\rangle
  -4\langle B^3\rangle\langle B\rangle
  -3\langle B^2\rangle^2
  \notag\\
  &\hspace{3.5cm}
  +12\langle B^2\rangle\langle B\rangle^2
  -6\langle B\rangle^4
  \bigr).
  \label{eq:qcd-baryon-fourth-cumulant}
\end{align}
If charge conjugation is an exact symmetry at \(x=0\) and sends
\(B\mapsto -B\), then all odd finite-volume cumulants vanish.
\end{proposition}

\begin{proof}
Because \(B_{a,V}\) is conserved, \(H_{a,V}\) and \(B_{a,V}\) can be
simultaneously decomposed by the spectral theorem, and differentiating
\eqref{eq:qcd-baryon-finite-x-partition} with respect to \(x\) is the
ordinary differentiation of the finite-regulator moment generating function
of the baryon charge in the thermal state.  For any observable \(Y\) commuting
with \(B\), in particular for \(Y=B^k\), one has
\begin{equation}
  \frac{\partial}{\partial x}\langle Y\rangle_x
  =
  \langle YB\rangle_x-\langle Y\rangle_x\langle B\rangle_x .
  \label{eq:qcd-baryon-source-derivative-expectation}
\end{equation}
Applying this identity to \(Y=1,B,B^2,B^3\), or equivalently differentiating
\(\log Z(x)\), gives the displayed cumulants.  If charge conjugation is an
exact symmetry at \(x=0\), the spectral measure of \(B\) in the thermal state
is invariant under \(B\mapsto -B\).  The moment generating function then
satisfies
\[
  Z(x)=Z(-x),
\]
so \(\log Z(x)\) is even near the origin and its odd derivatives vanish.
\end{proof}

\begin{definition}[Taylor expansion and critical diagnostic]
\label{def:qcd-baryon-taylor-critical-diagnostic}
When the thermodynamic continuum pressure is analytic at \(x=0\), write
\begin{equation}
  \frac{p(T,x)}{T^4}
  =
  \sum_{n=0}^{\infty}
  \frac{\chi_n^B(T)}{n!}x^n
  \label{eq:qcd-baryon-pressure-taylor}
\end{equation}
inside its disk of convergence.  The Cauchy--Hadamard radius of this germ is
\begin{equation}
  R(T)
  =
  \left[
  \limsup_{n\to\infty}
  \left|
  \frac{\chi_n^B(T)}{n!}
  \right|^{1/n}
  \right]^{-1}.
  \label{eq:qcd-baryon-cauchy-hadamard-radius}
\end{equation}
\end{definition}

\begin{proposition}[What a susceptibility radius establishes]
\label{prop:qcd-baryon-radius-diagnostic-status}
Assume that \(p(T,x)/T^4\) is the germ of a holomorphic function of complex
\(x\) at the origin and has Taylor coefficients
\(\chi_n^B(T)/n!\).  Then \(R(T)\) in
\eqref{eq:qcd-baryon-cauchy-hadamard-radius} is the radius of convergence of
the Taylor series.  The pressure is analytic for \(|x|<R(T)\).  If
independent input shows that the nearest limiting singularity is located at
a positive real point \(x_c\), then \(x_c=R(T)\) and this singularity is the
critical-point candidate seen by the Taylor expansion.  Without that
location input, the nearest singularity may lie at complex \(x\) and the
radius alone does not identify a real critical endpoint.
\end{proposition}

\begin{proof}
The first two assertions are the Cauchy--Hadamard theorem applied to the
power series \eqref{eq:qcd-baryon-pressure-taylor}.  A real critical point at
positive \(x_c\) must be a singularity of the same analytic germ after the
thermodynamic and continuum limits have been taken.  If it is independently
known to be the nearest singularity to the origin, its distance from the
origin is the convergence radius.  The final statement is the complementary
case: a power series records distance to the nearest complex obstruction,
whereas the physical critical endpoint is a statement about the real
\((T,\mu_B)\) plane and the limiting Gibbs/KMS states.
\end{proof}

\begin{remark}[Ratio estimators]
\label{rem:qcd-baryon-ratio-estimators}
For a charge-conjugation invariant theory the odd susceptibilities vanish at
\(x=0\), and one often forms even-order estimators
\[
  R_{2n}(T)
  =
  \sqrt{
    \frac{(2n+2)(2n+1)\chi_{2n}^B(T)}
         {\chi_{2n+2}^B(T)}
  }.
  \label{eq:qcd-baryon-ratio-estimator}
\]
If the even Taylor coefficients are eventually positive and have a single
dominant singular scale, these estimators can converge to the radius.  The
positivity and single-scale hypotheses are part of the diagnostic; they are
not consequences of the cumulant definitions alone.
\end{remark}

\section{Color Superconductivity and CFL Order}

At asymptotically large \(\mu_B\), quarks near the Fermi surface interact
through channels that can be attractive.  The BCS mechanism is then a
controlled weak-coupling instability of the Fermi surface, though the
precise gap prefactor is sensitive to magnetic gluon exchange and screening.

\begin{definition}[CFL condensate datum]
\label{def:qcd-cfl-condensate-datum}
For \(N_c=N_f=3\), a color-flavor-locked Higgs datum is a source-selected or
gauge-fixed representative of the gauge-covariant pairing tensor
\[
  \Phi^{ab}_{L,ij}
  =
  \langle q^a_{L i} C q^b_{L j}\rangle_J
  =
  \Delta_L\,\epsilon^{abA}\epsilon_{ijA},
  \qquad
  \Phi^{ab}_{R,ij}
  =
  \langle q^a_{R i} C q^b_{R j}\rangle_J
  =
  \Delta_R\,\epsilon^{abA}\epsilon_{ijA},
\]
with color indices \(a,b\), flavor indices \(i,j\), and spinor charge
conjugation matrix \(C\).  The subscript \(J\) means that an infinitesimal
pairing source and the infinite-volume limit are used before the source is
removed, exactly as in any spontaneous-order definition.  The tensor is not a
gauge-invariant local operator.  Its gauge-invariant content is the Higgs
pattern of the color gauge field, the residual global diagonal symmetry, and
the long-distance order of color-singlet composites built from the pairing
tensor.  The displayed tensor locks color and flavor: simultaneous color and
flavor rotations preserving the same diagonal \(SU(3)\) leave the
representative invariant.
\end{definition}

The color group is gauged, so broken color generators do not produce
physical Nambu--Goldstone particles; they Higgs the gauge bosons.  Ignoring
the anomalous axial \(U(1)\), the global symmetry breaking pattern is
\[
  SU(3)_L\times SU(3)_R\times U(1)_B
  \longrightarrow
  SU(3)_{\rm diagonal},
\]
and therefore gives
\[
  (8+8+1)-8=9
\]
physical Goldstone modes before quark-mass effects are included.  This
count is a group-theoretic statement about the order parameter.  Establishing
which region of QCD realizes it requires control of dense QCD, where direct
lattice importance sampling is obstructed by
Proposition~\ref{prop:qcd-finite-density-sign-problem}.

\section{Comparison of Confinement Criteria}

Several observables are often described by the same word ``confinement.''  In
this monograph they are kept distinct.

\begin{enumerate}
  \item The zero-temperature Wilson-loop area law is a statement about the
  large-loop asymptotics of a gauge-invariant loop operator and the static
  source potential.
  \item The thermal Polyakov-loop order parameter is a statement about
  center symmetry and the free energy of one static external charge.
  \item The 't Hooft loop probes magnetic or center-twist sectors and is
  dual, in favorable regimes, to Wilson-loop ordering.
  \item The center-vortex picture is a proposed infrared mechanism in which
  center-valued magnetic disorder controls the area law.
  \item The dual-superconductor picture is a proposed infrared mechanism in
  which condensation of magnetically charged objects squeezes electric flux.
\end{enumerate}

The first two items have clean operator definitions in lattice gauge theory.
The last three can be developed into sharp statements only after the
relevant line operators, disorder sectors, and continuum limits have been
specified.  Their value in this monograph is not as slogans, but as candidate
structures to be tested in examples and connected to the Wilsonian and BV
definitions of gauge theory developed earlier.

\begin{openproblem}[Continuum QCD phase diagram with controlled errors]
\label{op:qcd-phase-diagram-controlled-errors}
Construct the continuum QCD phase diagram in the \((T,\mu_B,m)\) variables
with mathematically controlled thermodynamic and continuum limits, including
the order of the pure-Yang--Mills deconfinement transition, the chiral
transition in massless limits, the physical crossover region, the existence
or nonexistence of a critical endpoint, and the color-superconducting
regions at high baryon density.
\end{openproblem}
